problem:  
Let \( G \) be a connected, simply connected **solvable** Lie group with Lie algebra \( \mathfrak{g} \).  
Denote by \( \mathcal{P}(\mathfrak{g}) \) the set of all left-invariant pseudo-Riemannian metrics \( g \) on \( G \) whose Levi-Civita connection satisfies \( \nabla \mathrm{Ric} = 0 \) but \( \mathrm{Ric} \neq \lambda g \) (so \( g \) is not Einstein).  

Let  
\[
\mathcal{M}(\mathfrak{g}) = \mathcal{P}(\mathfrak{g}) / \big( \mathbb{R}_{>0} \times \mathrm{Aut}(\mathfrak{g}) \big),
\]
where \(\mathbb{R}_{>0}\) acts by homothetic scaling and \(\mathrm{Aut}(\mathfrak{g})\) acts by pullback through Lie algebra automorphisms.

**Question:**  
Is the moduli space \(\mathcal{M}(\mathfrak{g})\) always discrete, or can there exist continuous families of inequivalent classes \([g_t]\) of left-invariant pseudo-Riemannian metrics with parallel Ricci tensor on a fixed solvable Lie group?  
If continuous families exist, can one characterize the algebraic features of \(\mathfrak{g}\) (such as the dimension of the center, presence of non-inner derivations, or degeneracies in the adjoint representation) that permit such deformations?  

Why is it a "good" problem:  
This problem probes a central but largely uncharted rigidity phenomenon in **homogeneous pseudo-Riemannian geometry**. For **Riemannian** solvmanifolds, Heber’s results show strong rigidity of Einstein metrics—moduli spaces are discrete, tightly controlled by the algebraic structure of \( \mathfrak{g} \). In contrast, the pseudo-Riemannian realm allows richer curvature behavior because signatures destroy many compactness arguments and balance conditions.  

The question asks whether similar rigidity (discrete moduli) persists under the weaker curvature constraint \( \nabla \mathrm{Ric} = 0 \), or whether solvable pseudo-Riemannian groups admit continuous one-parameter families of inequivalent Ricci-parallel metrics. Resolving this would clarify to what extent curvature restrictions in indefinite signatures control metric geometry and could reveal new moduli phenomena absent in the Riemannian setting.  

Its power lies in bridging:  
- **Lie algebraic structure** (derivation algebras, centers, cohomology of \(\mathfrak{g}\)),  
- **geometric analysis** (Ricci-parallel condition as an overdetermined PDE), and  
- **global pseudo-Riemannian geometry** (classification and moduli of homogeneous structures).  

Because it is both precisely formulated and conceptually profound—probing where the wall between rigidity and flexibility truly lies—it stands as a “good” and genuinely open research problem.